% !TeX root = ../main.tex
\begin{frame}{变形耦合构型}
  \begin{SlideCanvas}
    \LiteratureTimeAxis
    \LiteratureTimeNode{255}{2018}
    \LiteratureTimeNode{705}{2020}
    \PointAt{55}{153}{400}{Wu等：反四手性}
    \PointAt{505}{153}{400}{Ruan等：手性约束}
    \CenteredPicAt{255}{278}{390}{188}{assets/wu2018_fig1.jpg}
    \CenteredPicAt{705}{278}{390}{188}{assets/chiral2020_fig1.jpg}
    \DetailAt{55}{386}{400}{提出圆形／椭圆节点的反四手性及层级构型，研究其力学响应。}
    \DetailAt{505}{386}{400}{通过周向约束杆调节变形耦合，并用狭窄血管模型评价部署表现。}
    \Sources{\RefWu，图1；\RefRuan，图1。叙述依据期刊摘要与原图。\\时间轴按卷期年份排序，节点间距仅作示意。}
  \end{SlideCanvas}
  \BICConclusionBand{手性构型与约束设计提供了\BICBandEmph{调节周向—轴向耦合的基础}}
  \note{按2018到2020的时间轴引出两篇研究，每篇各配一张原文图。Wu的图1为带椭圆节点的反四手性单元，论文摘要还提出层级构型；不把单元图标成已验证的临床器件。Ruan图1展示P1、P2、PS1及其代表单元，说明手性结构与周向约束的组合。两篇的力学与部署描述限于公开摘要和图件，不补写未核对的临床获益。年份按卷期，节点间距不表示真实时间跨度。}
\end{frame}
